\section{Introduction}

Anthropogenic climate change, driven by greenhouse gas (GHG) emissions, constitutes an economic externality due to its significant economic damages.
Integrated Assessment Models (IAMs) provide frameworks that model the interactions between economic policy and climate, helping to develop policies that mitigate these externalities. However, as IAMs increase in complexity, they can encounter computational challenges commonly referred to as the \enquote{curse of dimensionality} \parencite{Bellman1961}. With traditional solution methods, current IAMs often rely on restrictive assumptions to remain computationally tractable. This thesis aims to address these challenges by employing state-of-the-art deep learning methods that could potentially allow for more complex IAMs, reducing the need for restrictive assumptions.

One of the most notable applications of IAMs is in the assessment reports of the Intergovernmental Panel on Climate Change (IPCC), which evaluate the state of the global climate and inform policy decisions \parencite{IPCC2023}. IAMs have also been used to calculate the official social cost of carbon (SCC) for the United States \parencite{IWGSCC2013}, further underscoring their significance in economic policy formulation. 

However, despite their extensive use, these models can face substantial computational challenges. As \textcite{Azinovic2022DEQN} notes, these challenges arise when the system is stochastic, the state space is high-dimensional, there are nonlinearities in the equilibrium functions and complex distributions of state variables within the model's state space (ergodic set of states). Advancements in deep learning have begun to address these challenges, particularly within the realm of reinforcement learning (RL). These methods have been recently applied to and developed for economic models \parencite{Maliar2021, Azinovic2022DEQN, Valaitis2024}.

This thesis extends the application of Deep Equilibrium Nets (DEQNs), as presented in \textcite{Azinovic2022DEQN}, by applying them to compute recursive economic equilibrium within the ACE model in a discrete-time setting. To the best of my knowledge, this work joins a narrow field of research using DEQNs for IAMs, where the only previously published papers are the ones by \textcite{Folini2024ClimateInCE, friedl2023_deepUQ}. 

Distinct from previous efforts that predominantly focused on DICE models, this thesis applies DEQNs to the Analytic Climate Economy (ACE) model, as presented in \textcite{TraegerChristian2023ACE}. The advantage of applying this method to ACE is that it gives an analytic solution that can be used to analyze the quality of the numeric solution. Additionally, ACE is general in its structure, making my derivations easily extendable to other IAMs as well. 

Furthermore, this thesis emphasizes the relationship between DEQNs and deep RL, a connection that has not been extensively explored in existing literature. This linkage is critical as it opens the door to integrating more advanced algorithms from the RL field, potentially improving existing methodologies for computing economic equilibria. Accompanying the thesis is a publicly available GitHub repository that serves as a comprehensive resource for future research in this area\footnote{The GitHub repository is available here: \url{https://github.com/hermanvest/DNN-for-ACE}}.

This thesis identifies crucial areas for improvement in the solution method, particularly emphasizing the need for more robust normalization techniques in the pre- and post-processing of state and control variables. Proper application of these techniques is important for ensuring effective learning. Improved normalization could facilitate faster model convergence and result in more stable training processes.

The thesis is structured as follows: Section two provides a comprehensive review of the literature on IAMs related to ACE and explores computational methods for estimating economic equilibria, with a specific emphasis on RL. Section three delves into preliminary knowledge, covering RL concepts necessary to understand DEQNs. In section four, the general structure of the ACE model is discussed. Section five presents the necessary theoretical derivations required for mapping the general structure of ACE onto DEQNs. Section six discusses the results and findings. The thesis concludes in section seven. For detailed derivations of the specific models, supplementary figures and simulation initialization, please refer to the appendix of the thesis.